% Recuperación aditiva del propietario N57, con enlace a las pruebas
% de cilindros y firma que ya están incluidas en el manuscrito.
\subsection{Cociclo de acarreo y levantamiento entero de la firma regional}
\label{carry27:seccion}

La firma regional conserva simultáneamente el residuo y el cociente.
La siguiente construcción hace explícita la ley algebraica de este
último dato. APP proporciona las evaluaciones enteras sobre sus dos
hojas; TRIT conserva orientación y régimen; el transporte TPK
produce las bandas regionales y actualiza su memoria. En una
columna de esas bandas, la división euclídea determina el acarreo.
La construcción utiliza los representantes \(\{0,1,2\}\) de
\(\mathbb F_3\); la representación orientada del TRIT conserva
por separado su signo y su cociente de transporte.

Sean \((p,e,f)\in\{0,1,2\}^3\) las entradas de una columna
regional. En este apartado, la letra \(e\) designa la entrada
ternaria de la región de propagación. Definimos
\begin{equation}
\begin{aligned}
q&=[p+e+f]_3,& c&=(p+e+f-q)/3,\\
a&=[p+e-f]_3,& h&=(p+e-f-a)/3,
\end{aligned}
\label{carry27:division}
\end{equation}
donde \([z]_3\in\{0,1,2\}\) es el resto de cualquier entero \(z\).
La variable \(a\) es la coordenada axial de esta columna. La
constante \(\alpha_{\rm HMT}\) conserva su construcción
dodecafásica posterior.

\begin{proposition}[Composición local de los acarreos]
\label{carry27:cociclo}
Para \(x,y,z\in\{0,1,2\}\), sean
\[
\kappa_3(x,y)=\frac{x+y-[x+y]_3}{3},\qquad
\beta_3(x,y)=\frac{x-y-[x-y]_3}{3}.
\]
Con \(r=[p+e]_3\), las dos leyes son
\[
c=\kappa_3(p,e)+\kappa_3(r,f),\qquad
h=\kappa_3(p,e)+\beta_3(r,f).
\]
Además,
\begin{equation}
\kappa_3(x,y)+\kappa_3([x+y]_3,z)
=\kappa_3(y,z)+\kappa_3(x,[y+z]_3).
\label{carry27:identidad}
\end{equation}
\end{proposition}

\begin{proof}
La igualdad \(p+e=r+3\kappa_3(p,e)\), sumada a
\(r+f=q+3\kappa_3(r,f)\), da la primera fórmula.
Restando \(f\) y usando
\(r-f=a+3\beta_3(r,f)\) se obtiene la segunda.
Ambos miembros de \eqref{carry27:identidad} son
\((x+y+z-[x+y+z]_3)/3\). Por tanto, la composición
conserva el mismo cociente cualquiera que sea la asociación
de la suma.
\end{proof}

Si \(s_3:\mathbb F_3\to\mathbb Z\) es la sección de
representantes, entonces
\[
s_3(x)+s_3(y)-s_3(x+y)=3\kappa_3(x,y).
\]
Éste es el cociclo de la extensión
\(0\to\mathbb Z\xrightarrow{\times3}\mathbb Z\to\mathbb F_3\to0\).
Su aparición procede del levantamiento de una operación residual
a sus valores enteros. La coordenada visible y la memoria del
múltiplo de tres quedan relacionadas por una identidad exacta.

\Needspace{10\baselineskip} % S27_LAYOUT_TABLE
\begin{longtable}{ccc|rrrr}
\caption{Ley completa de las \(27\) columnas regionales.
Los valores de \(h\) se muestran como enteros.}
\label{carry27:tabla}\\
\toprule
\(p\)&\(e\)&\(f\)&\(q\)&\(c\)&\(a\)&\(h\)\\
\midrule\endfirsthead
\toprule
\(p\)&\(e\)&\(f\)&\(q\)&\(c\)&\(a\)&\(h\)\\
\midrule\endhead
\bottomrule\endfoot
0&0&0&0&0&0&0\\
0&0&1&1&0&2&-1\\
0&0&2&2&0&1&-1\\
0&1&0&1&0&1&0\\
0&1&1&2&0&0&0\\
0&1&2&0&1&2&-1\\
0&2&0&2&0&2&0\\
0&2&1&0&1&1&0\\
0&2&2&1&1&0&0\\
1&0&0&1&0&1&0\\
1&0&1&2&0&0&0\\
1&0&2&0&1&2&-1\\
1&1&0&2&0&2&0\\
1&1&1&0&1&1&0\\
1&1&2&1&1&0&0\\
1&2&0&0&1&0&1\\
1&2&1&1&1&2&0\\
1&2&2&2&1&1&0\\
2&0&0&2&0&2&0\\
2&0&1&0&1&1&0\\
2&0&2&1&1&0&0\\
2&1&0&0&1&0&1\\
2&1&1&1&1&2&0\\
2&1&2&2&1&1&0\\
2&2&0&1&1&1&1\\
2&2&1&2&1&0&1\\
2&2&2&0&2&2&0\\
\end{longtable}

\begin{theorem}[Componente algebraica adicional del acarreo]
\label{carry27:rango}
Sobre \(\mathbb F_3\), las funciones \(q,a\) son lineales en
\((p,e,f)\). Las reducciones de \(c,h\) tienen las siguientes
representaciones polinómicas:
\begin{align}
\bar c={}&2ef+2ef^2+2e^2f+2pf+2pf^2\notag\\
 &+2pe+pef+2pe^2+2p^2f+2p^2e,
\label{carry27:polinomio-c}\\
\bar h={}&2f^2+ef+2ef^2+e^2f+pf+2pf^2\notag\\
 &+2pe+2pef+2pe^2+p^2f+2p^2e.
\label{carry27:polinomio-h}
\end{align}
Si \(V\) es la matriz de evaluación de \(1,p,e,f\) en las
\(27\) columnas, se cumple
\[
\operatorname{rank}_{\mathbb F_3}V=4,\qquad
\operatorname{rank}_{\mathbb F_3}(V\mid\bar c)=5,\qquad
\operatorname{rank}_{\mathbb F_3}(V\mid\bar h)=5.
\]
\end{theorem}

\begin{proof}
La sustitución de las \(27\) ternas de la tabla en los dos
polinomios reproduce \(\bar c,\bar h\). La representación con
grado a lo sumo dos en cada variable es única: un polinomio de
grado a lo sumo dos en una variable que se anula en los tres
elementos de \(\mathbb F_3\) es nulo; aplicando sucesivamente
este argumento a las tres variables se obtiene la independencia
de los \(27\) monomios.

Las evaluaciones en \(0,(1,0,0),(0,1,0),(0,0,1)\) prueban
la independencia de \(1,p,e,f\). La función \(\bar c\) vale
cero en esas cuatro entradas y vale uno en \((1,1,1)\);
por ello queda fuera de su espacio afín de evaluación.
Para \(\bar h\), las tres primeras entradas fuerzan los
coeficientes constante, de \(p\) y de \(e\) a cero, y
\(\bar h(0,0,1)=2\) fuerza el coeficiente de \(f\) a dos.
Esa función afín valdría uno en \((0,0,2)\), mientras
\(\bar h(0,0,2)=2\). Cada columna adicional aumenta,
por tanto, el rango en uno.
\end{proof}

La prueba especifica qué información aporta el cociente. La
superposición residual produce \(q,a\); la evaluación entera
incorpora además \(c,h\). Para \(c\in\{0,1,2\}\) su residuo
determina el valor entero, y para \(h\in\{-1,0,1\}\) lo
determina junto con esta elección declarada de representantes.
Los polinomios describen las reducciones; la memoria íntegra
conserva también esos levantamientos.

\subsubsection{Transporte de los cocientes regionales}

Sea \(M\) una matriz entera de transporte ya fijada en una etapa
del desarrollo, y sean \(x_t^\chi\) sus estados fila para
\(\chi\in\{\pi,e,\varphi\}\). Para cada coordenada se realiza
la división
\[
y_t^\chi=x_t^\chi M,\qquad
u_t^\chi=[y_t^\chi]_3,\qquad
x_{t+1}^\chi=(y_t^\chi-u_t^\chi)/3.
\]
Definamos \(X_t^+=x_t^\pi+x_t^e+x_t^\varphi\) y
\(X_t^\sigma=x_t^\pi+x_t^e-x_t^\varphi\).
Las leyes anteriores, aplicadas columna a columna, dan
\begin{equation}
\begin{aligned}
X_{t+1}^+&=\frac{X_t^+M-q_t}{3}-c_t,\\
X_{t+1}^\sigma&=\frac{X_t^\sigma M-a_t}{3}-h_t.
\end{aligned}
\label{carry27:transporte}
\end{equation}
En efecto, la suma de los tres residuos es \(q_t+3c_t\);
la combinación firmada es \(a_t+3h_t\). Sustituir ambas
identidades en la división de los estados demuestra
\eqref{carry27:transporte}. El enunciado tiene por dominio
los transportes enteros declarados. Las emisiones regionales
y su continuidad conservan los lectores coinductivos
previamente construidos.

Como evaluación completa en una banda de seis columnas,
\[
u^\pi=021020,\quad u^e=001220,\quad u^\varphi=100210
\]
producen \(q=122120\), \(c=000110\), \(a=222000\) y
\(h=(-1,0,0,0,1,0)\). La orientación
\(\sigma=(1,1,-1)\) posee residuo doble
\(2\sigma=(2,2,1)\) en \(\mathbb F_3^3\). Esta última
igualdad registra el convenio de signo y se distingue de las
bandas de seis columnas.

\subsubsection{Memoria de frontera y lectura del intervalo}

La ley local de acarreo se acopla al refinamiento cilíndrico
ya demostrado. Un ejemplo muestra la función del intervalo
completo:
\[
\left[\frac1{729},\frac2{729}\right)
\cap
\left[\frac2{1000},\frac3{1000}\right)
=
\left[\frac2{1000},\frac2{729}\right)\ne\varnothing.
\]
En cambio,
\(\lfloor1000/729\rfloor=1\).
La lectura del extremo inferior devuelve el bloque \(1\),
mientras la compatibilidad intervalar admite también la
celda \(2\). El estado de frontera conserva las dos cotas
necesarias para decidir esa compatibilidad.

Para una columna, fijar \(q\) y \(a\) fija
\[
f=2(q-a),\qquad p+e=q-f\quad\text{en }\mathbb F_3.
\]
Cada una de las nueve fibras de \((q,a)\) contiene exactamente
tres ternas: elegir \(p\) determina \(e\) y deja \(f\) fijo.
En seis columnas, antes de aplicar restricciones adicionales,
la libertad es el producto de esas seis fibras y tiene
cardinal \(3^6=729\). La fijación se refiere a alternativas
de una misma frontera. Al avanzar en profundidad cambian
el registro y su memoria según el transporte. De esta
forma, el acarreo local, las cotas del cilindro y la
incidencia regional actúan conjuntamente sobre la misma
historia aritmética.
